% Abridged version, Section 7, part 1 of 2: formal stationary phase, the full-Stirling series, Theorem P(a). Continued in part 2.
\section{The connection problem}
\label{sec:con}

The canonical gauge is normalised at the conifold point $u=0$, the rows of the formal quintic at $q=0$; the normalisation series $\Dn_0$ of \cref{thm:gau-R0} measures the difference. This section reduces $\Dn_0=1$ to one identity at $q=0$. The bridge is the saddle series $c(z;L)$ of a Mellin--Barnes summand, with hatted form $\PhiSP$, which is defined for all $q$ at once: Theorem~P shows that $\PhiSP$ is $\Phican$, and the glue gives $\Dn_0(z)=\PhiSP(z,1)/h(z)$. We work in colour~$0$. This section follows \Lr{sec:con}.

%-------------------------------------------------------------------------------
\subsection{Formal stationary phase}
\label{sec:con-fsp}

Stationary phase is used as an algebraic operation on Taylor series \Lr{sec:fsp}; nothing is evaluated at $z\neq0$. Let $A$ be a commutative $\bbQ$-algebra and $c\in A^n$. A \emph{phase} $F\in A\fps{x-c}$ has the \emph{non-degenerate critical point} $c$ if $\nabla F(c)=0$ and $\mathsf H:=\Hess F(c)\in GL_n(A)$. Write $F(c+t)=F(c)+\frac12t^{\mathrm T}\mathsf Ht+F_{\geq3}(t)$ and $C:=-\mathsf H^{-1}$. For an amplitude $a\in A\fps{x-c}\lser z$ the \emph{normalised stationary phase} is
\[
\SPn_c[F;a]:=E_{zC}\bigl[\exp\bigl(F_{\geq3}(t)/z\bigr)\,a(c+t)\bigr]\in A\lser z,\qquad E_{zC}[t^\alpha z^j]:=\langle t^\alpha\rangle_C\,z^{j+|\alpha|/2},
\]
where $\langle t^\alpha\rangle_C$ is the Gaussian moment with covariance $C$, the sum over the perfect matchings of the products of the entries of $C$ \Lr{def:fsp-ring,def:fsp-wick}. If $a=\sum_{m\geq0}a_mz^m$, its constant term is $a_0(c)$; if this is a unit, $\SPh_c[F;a]:=\SPn_c[F;a]/a_0(c)\in1+zA\fps z$.  The coefficients of $\SPn_c[F;a]$ are universal polynomials over $\bbQ$ in the entries of $C$ and in finitely many Taylor coefficients of $F$, of order at least $3$, and of $a$ \Lr{lem:fsp-univ}; in particular they do not see constants added to $F$. By \Lr{sec:fsp}, the stationary phase vanishes on exact forms $\dd(\re^{F/z}\iota)$; it is invariant under formal changes of variables and under translations by multiples of $z$; it can be iterated variable by variable (a formal Fubini theorem); it commutes with derivatives in a parameter. The \emph{formal Stirling form} is the expression $\Gst(Y):=\exp\bigl((Y-\frac12)\log Y-Y+\frac12\log2\pi+\sum_{k\geq1}\bk_kY^{1-2k}\bigr)$, with $\Gst(1+w):=w\,\Gst(w)$; it satisfies the functional equation, and it is the stationary phase of Euler's integral: for a unit $a$ and an integer $M$,
\[
\SPn_a\bigl[a\log x-x;\ x^{M-1}\bigr]=a^{M-1}\,\mathrm{Poch}_M(z/a)\,\exp\Bigl(\sum_{k\geq1}\bk_k\,(z/a)^{2k-1}\Bigr)
\]
\Lr{lem:fsp-stirfe,lem:fsp-stirling}, where $\mathrm{Poch}_M(w)$ is the polynomial or rational function with $\mathrm{Poch}_{m+1}(w)=(1+mw)\mathrm{Poch}_m(w)$ and $\mathrm{Poch}_0=\mathrm{Poch}_1=1$.

%-------------------------------------------------------------------------------
\subsection{The full-Stirling saddle series and \texorpdfstring{$\PhiSP$}{the full-Stirling WKB function}}
\label{sec:con-form}

\begin{defn}[The full-Stirling saddle series; {\Lr{def:gam-c,prop:gam-rational}}]
\label{def:gam-c}
Let $\Cstar:=\bbC\setminus\{0,1,\zeta,\zeta^2,\zeta^3,\zeta^4\}$ and $q(L):=(1-L^{-5})/5^5$. For $L\in\Cstar$ fix an open disc $\bbD_L$ around $L$ containing neither $0$ nor any $\zeta^j$, holomorphic branches of $\log(5\sigma)$ and of $\log(\zeta^j-\sigma)$ on $\bbD_L$, and put $\lambda:=-5\log(5L)+\sum_{j=0}^4\log(\zeta^j-L)$, so that $\re^\lambda=-q(L)$. On $\bbD_L$ define
\begin{align*}
\mathsf S_0(\sigma)&:=\lambda\sigma+5\sigma\log(5\sigma)-5\sigma+\sum_{j=0}^4\bigl[(\zeta^j-\sigma)\log(\zeta^j-\sigma)-(\zeta^j-\sigma)\bigr],\\
\mathsf S_1(\sigma)&:=\tfrac12\log(5\sigma)-\tfrac12\sum_{j=0}^4\log(\zeta^j-\sigma),\qquad\mathsf S_{2k}(\sigma):=\bk_k\Bigl[(5\sigma)^{1-2k}+\sum_{j=0}^4(\zeta^j-\sigma)^{1-2k}\Bigr]\quad(k\geq1),
\end{align*}
and $A_{\mathrm{St}}(\sigma,z):=\exp\bigl(\mathsf S_1(\sigma)+\sum_{k\geq1}z^{2k-1}\mathsf S_{2k}(\sigma)\bigr)$.
Then $\mathsf S_0'(L)=0$ and $\mathsf S_0''=5/(\sigma(1-\sigma^5))$, which does not vanish at $L$. The \emph{full-Stirling saddle series} is
\[
c(z;L):=\SPh_L[\mathsf S_0;A_{\mathrm{St}}](z)=\sum_{n\geq0}c_n(L)z^n,\qquad c_0(L)=1 .
\]
It does not depend on the disc nor on the branches of the logarithms, and every $c_n$ is the value at $L$ of an element of $\AL=\bbQ[L,L^{-1},(L^5-1)^{-1}]$.
\end{defn}

Up to a factor $(2\pi)^3z^2$, $\exp(\mathsf S_0/z)A_{\mathrm{St}}$ is the product $\re^{\lambda\sigma/z}\,\Gst(1+5\sigma/z)\prod_j\Gst((\zeta^j-\sigma)/z)$: this is the summand of \cref{thm:int-gamma} with all six Gamma factors replaced by $\Gst$. 

By the reflection formula for $\Gst$, the same series is the saddle series of the colour-$0$ summand $q^{\sigma/z}\,\Gamma(1+5\sigma/z)/\prod_j\Gamma(1+(\sigma-\zeta^j)/z)$ \Lr{lem:con-form}~(d). In the variable $x=1/\sigma$ this summand, with its Stirling expansions, is $\re^{-\cW/z}\fa\,\dd x$ up to a factor independent of $x$, where, with $\Xi=-\log(1-u)$,
\begin{gather*}
\cW(x):=\frac{\Xi}x+\sum_{m\geq1}\frac{x^{5m-1}}{m(5m-1)},\\
\log\fa(x,z)=-\tfrac12\log(1-x^5)+\sum_{k\geq1}\bk_k(xz)^{2k-1}\Bigl[5^{1-2k}-5\sum_{P\geq0}\binom{2k-2+5P}{5P}x^{5P}\Bigr].
\end{gather*}
In the hatted variables $x=sX$, $z=s^4\zh$ one has $\cW/z=\hat\cW/\zh$ and $\hat\fa(X;\zh):=\fa(sX,s^4\zh)\in\bbQ[X,\zh]\fps u$, with
\[
\hat\cW(X;u)=\sum_{m\geq1}\frac{u^{m-1}}m\Bigl[\frac1X+\frac{X^{5m-1}}{5m-1}\Bigr],\qquad\hat\cW'(1)=0,\qquad\hat\cW''(1)=\frac5{1-u},\qquad\hat\cW(1)=\frac5{4\,\tau(u)} :
\]
the point $X=1$ is critical for every $u$, and by \cref{lem:set-canon} the critical exponential is $\re^{-\hat\cW(1)/\zh}=\re^{(u_0-u_0^{\mathrm c})/z}$ \Lr{lem:con-form}.

\begin{defn}[The full-Stirling WKB function and the residue; {\Lr{def:con-PhiSP}}]
\label{def:con-PhiSP}
The \emph{full-Stirling WKB function} is the normalised stationary-phase series, with constant term $1$, of $\re^{-\hat\cW/\zh}\hat\fa\,\dd X$ at $X=1$:
\[
\PhiSP(\zh,u):=\SPh_{X=1}[-\hat\cW;\hat\fa](\zh)=\sum_{n\geq0}\mathrm{ph}^{\mathrm{SP}}_n(u)\zh^n,\qquad\mathrm{ph}^{\mathrm{SP}}_0=1,\quad\mathrm{ph}^{\mathrm{SP}}_n\in\bbQ\fps u .
\]
The \emph{residue} of the flat $1$-form is $\cV:=\Res_{x=0}\re^{-\cW/z}\fa\,\dd x:=\sum_{k\geq1}\frac{(-\Xi/z)^k}{k!}\,[x^{k-1}]\bigl(\re^{-\cG/z}\fa\bigr)\in z^{-1}\bbQ[z^{-1}]\fps u$, where $\cG:=\cW-\Xi/x$.
\end{defn}

So one $1$-form carries two functionals, the stationary phase at the critical point and the residue at the pole; for the model form $p^m\re^{-\Wv/z}\dd p$ these are the model rows and the vanishing solution $\IC_m$.

%-------------------------------------------------------------------------------
\subsection{Theorem P}
\label{sec:con-P}

The ring $\AL$ carries the derivation $\Dq$ with $\Dq L=(L^6-L)/5$, for which $\Dq q=q$; let $\jmath\colon\AL\to\bbQ\lser s$ be the injective ring homomorphism with $\jmath(L)=s^{-1}$, and write $\PFfive^L$ for $\PFfive$ with $z\Dq$ replaced by $z\Dq+L$.

\begin{thm}[Theorem P(a); {\Lr{thm:con-Pa}}]\genusstatus{AI-proposed, unchecked}
\label{thm:con-Pa}
Let $c(z;L)=\sum_nc_n(L)z^n$ be the full-Stirling saddle series, with $c_n\in\AL$.
\begin{enumerate}
\item $\PFfive^{L}\bigl(L\,c(z;L)\bigr)=0$ in $\AL\fps z$.
\item In $\bbQ\lser s\fps z$ one has $\jmath(c)(z)=\PhiSP(s^{-4}z,s^5)$; that is, $\jmath(c_n)=s^{-4n}\mathrm{ph}^{\mathrm{SP}}_n$ for every $n$.
\item $\PhiSP$ solves HPF, \eqref{eq:gau-HPF}, in $\bbQ\fps u\fps{\zh}$.
\item Every $\mathrm{ph}^{\mathrm{SP}}_n$ is a polynomial in $u$; hence $c_n(L)=L^{4n}\mathrm{ph}^{\mathrm{SP}}_n(L^{-5})$ in $\AL$, and $c_n\in\bbQ[L,L^{-1}]$.
\end{enumerate}
The proof uses no input from outside [L].
\end{thm}

\begin{sketch}
(i) The series $c$ depends on $q$ only through the factor $(-q)^{\sigma/z}$ of the summand, so $z\Dq$ acts under the stationary phase as multiplication by $\sigma$: for $c\langle\chi\rangle:=\SPn_L[\mathsf S_0;\chi A_{\mathrm{St}}/A_{\mathrm{St}}(L,0)]$, $\chi\in\bbQ[\sigma,z]$, one has $(z\Dq+L)(L\,c\langle\chi\rangle)=L\,c\langle\sigma\chi\rangle$ (the derivative in a parameter, with the moving critical point $L$). Hence $\PFfive^L(L\,c)=L\bigl(c\langle\sigma^5-1\rangle-c\langle q\prod_{k=1}^5(5\sigma+kz)\rangle\bigr)$. The two terms are equal, because the translation $\sigma\mapsto\sigma+z$ leaves the stationary phase invariant and, by the functional equation of $\Gst$, turns the amplitude $(\sigma^5-1)A_{\mathrm{St}}$ into $q\prod_k(5\sigma+kz)A_{\mathrm{St}}$. (ii) is the change of variables $\sigma=L/X$, under which the phase $\mathsf S_0$ becomes $-s^4(\hat\cW-\hat\cW(1))$ and the normalised amplitude becomes $\hat\fa/\hat\fa(1;0)$. (iii) follows from (i) and (ii) by the chain rule of \cref{lem:gau-HPF}, and (iv) from \cref{lem:gau-HPF}~(iii).
\end{sketch}

% Abridged version, Section 7, part 2 of 2: Theorem P(b) and the glue.

Theorem~P(b) identifies $\PhiSP$ with $\Phican$. Both solve HPF with constant term $1$ and coefficients in $\bbQ\fps u$, so by \cref{lem:gau-HPF}~(iii) they differ by a series $\mathsf s(z)\in1+z^5\bbQ\fps{z^5}$, which has to be shown to be $1$. The equation does not determine it; it is determined by a second functional on the same $1$-form, the residue at the pole, compared with the gauge identity (G).

\begin{lem}[The uniformising map and the reduction; {\Lr{lem:con-cfu,lem:con-red}}, abridged]\genusstatus{AI-proposed, unchecked}
\label{lem:con-cfu}
There are unique $Y_k\in\bbQ[T]$ ($k\geq1$) and $V\in1+u\bbQ\fps u$ such that $Y:=X\bigl(1+\sum_{k\geq1}u^kY_k(X^5)\bigr)$ satisfies
\[
\hat\cW(X;u)=V/Y+Y^4/4=W_V(Y)\qquad\text{in }\bbQ[X,X^{-1}]\fps u.
\]
Moreover $V=\tau^{-5/4}$, and the point $X=1$ goes to the critical point $Y_c:=Y(1)=V^{1/5}=\tau^{-1/4}$ of $W_V$. Let $\fa_Y(Y;\zh):=\hat\fa(X(Y);\zh)\,X'(Y)\in\bbQ[Y,\zh]\fps u$ be the hatted amplitude in the variable $Y$. There are $\pi_m\in\bbQ[\zh]\fps u$ ($m=-1,\dots,3$) and $\mathcal E\in Y\bbQ[Y,\zh]\fps u$ with
\[
\fa_Y\,\re^{-W_V/\zh}\,\dd Y=\sum_{m=-1}^{3}\pi_m\,Y^m\re^{-W_V/\zh}\,\dd Y+\dd\bigl(\mathcal E\,\re^{-W_V/\zh}\bigr),
\]
and $\pi_m$ contains only powers $\zh^k$ with $k\equiv m\bmod5$.
\end{lem}

Here $V$ is identified by the critical value $\hat\cW(1)=\frac54\tau^{-1}=\frac54V^{4/5}$; in the flat variables $p=sY$ the parameter of the model potential is again $uV=v(u)$. The reduction uses (Mod2), in the form of the exact forms $\zh\,\dd(Y^{n-3}\re^{-W_V/\zh})$, to lower the degree in $Y$ to the five exponents $-1,\dots,3$, which are the charges $\nu_b$.

\begin{lem}[The two functionals on the reduced form; {\Lr{lem:con-S2,lem:con-id,prop:con-res}}, abridged]\genusstatus{AI-proposed, unchecked}
\label{lem:con-id}
Work in a conifold frame.
\begin{enumerate}
\item (The model identity.) For $W_1(\tilde p):=\tilde p^4/4+1/\tilde p$ and every $b\in\bbZ/5$, $\SPn_{\tilde p=1}[-W_1;\tilde p^{\nu_b}](t)=\varepsilon_b\Fs_b(t)$.
\item $\PhiSP=\sum_b\varepsilon_b\,\pi_{\nu_b}\,\tau^{(1-\nu_b)/4}\Fs_b(\tau\zh)$.
\item Let $P$ be a graded regular gauge with coefficients in $\bbQ\fps u$ from the pulled-back model to the quintic system, and let $\Phi_P$ be $\varpi_0$ times the colour-$0$ entry of the hatted row $0$ of $P\,\Qr^{\mathrm m}(v(u))$; for $P=\Pc$ this is $\Phican$. Then $\Phi_P/\varpi_0=\sum_b\tau^{(1-\nu_b)/4}\,G_b\,\Fs_b(\tau\zh)$ with $G_b:=s^{\nu_b}P_{0b}(u,s^4\zh)$, and such an expansion, with coefficients that contain only powers $\zh^k$, $k\equiv-b\bmod5$, is unique. Hence, if $\Phi_P=\PhiSP$, then $G_b=\varepsilon_b\pi_{\nu_b}/\varpi_0$ for every $b$.
\item (The residue.) If $P$ is as in (iii) and $\Phi_P=\PhiSP$, then $\cV=\varpi_0\,\bigl(P\,\YC(v(u),z)\bigr)_0$ in $\bbQ\lser z\fps u$.
\end{enumerate}
\end{lem}

\begin{sketch}
(i) The stationary phases $J_m$ of $p^m\re^{-\Wv/z}\dd p$ at the critical point $p=\varsigma$ satisfy (Mod1), by the derivative in the parameter $v$, and (Mod2), by the vanishing on exact forms; so they solve the model system, and the uniqueness of the T-system (\cref{prop:set-modrows}) identifies them. (ii) The stationary phase is invariant under the uniformising map and vanishes on the exact part, so \cref{lem:con-cfu} and (i), rescaled to $Y_c$, give the formula. (iv) The formal residue at the pole is also invariant under the uniformising map and vanishes on exact forms; so the residue of the hatted form is $\sum_m\pi_m\Res_Y[Y^m\re^{-W_V/\zh}\dd Y]$, and the model residues are the $\IC_m(v(u),z)$ up to powers of $s$. The same coefficients $\pi_{\nu_b}$, which are the entries of $P$ by (iii), now multiply the entries of the vanishing solution: the sum is $\varpi_0\sum_bP_{0b}\YC_b$.
\end{sketch}

\begin{thm}[Theorem P(b); {\Lr{thm:con-Pb}}]\genusstatus{AI-proposed, unchecked}
\label{thm:con-Pb}
In every conifold frame of \cref{def:set-frames}, $\PhiSP=\Phican$ in $\bbQ\fps u\fps{\zh}$. The proof does not use (QRR).
\end{thm}

\begin{proof}[Proof, given the lemmas]
Let $\PhiSP=\Phican\,\mathsf s(z)$ with $\mathsf s\in1+z^5\bbQ\fps{z^5}$, and put $P^{\mathrm{SP}}:=\mathsf s(z)\Pc$. By \cref{thm:gau-U} it is a graded regular gauge, and $\Phi_{P^{\mathrm{SP}}}=\mathsf s(z)\Phican=\PhiSP$. The gauge identity $\Pc\YC(v(u),z)=\YCQ$ (\cref{thm:gau-G}), multiplied by $\mathsf s(z)$, gives $(P^{\mathrm{SP}}\YC)_0=\mathsf s(z)(\YCQ)_0$, and \cref{lem:con-id}~(iv) for $P=P^{\mathrm{SP}}$ gives
\[
\cV=\varpi_0\,\mathsf s(z)\,(\YCQ)_0\qquad\text{in }\bbQ\lser z\fps u .
\]
Compare the coefficients of $u^1$. On the left only $k=1$ contributes in \cref{def:con-PhiSP}, because $\Xi=u+O(u^2)$, and $[x^0](\re^{-\cG/z}\fa)=1$; so $[u^1]\cV=-1/z$. On the right, $[u^0](\YCQ)_0=0$ and $[u^1](\YCQ)_0=-1/z$, since only $y^{(1)}_0=-\tc$ contributes at $u^1$ and $\tc=u+O(u^2)$ (\cref{lem:gau-vanq}, \Lr{lem:gau-vanq}); with $\varpi_0(0)=1$ the right side has the $u^1$-coefficient $-\mathsf s(z)/z$. Hence $\mathsf s=1$.
\end{proof}

So the freedom $\cn(z^5)$ left by the differential equation is fixed on both sides by the same normalisation, the residue at the pole: on the side of the gauge by (G), on the side of the Mellin--Barnes summand by the explicit residue $\cV$.

%-------------------------------------------------------------------------------
\subsection{The glue}
\label{sec:con-glue}

For a colour-$0$ WKB function with polynomial coefficients write $\Phi(z,1):=\sum_n\mathrm{ph}_n(1)z^n$, its value at $u=1$, where $\zh=z$.

\begin{thm}[The glue; {\Lr{thm:con-glue}}]\genusstatus{AI-proposed, unchecked}
\label{thm:con-glue}
In $\bbQ(\zeta)\fps z$, $h(z)\,\Dn_0(z)=\Phican(z,1)$; hence $\Dn_0\in1+z\bbQ\fps z$, and, by \cref{thm:con-Pb}, $\Dn_0(z)=\PhiSP(z,1)/h(z)$.
\end{thm}

\begin{proof}
By \cref{thm:gau-R0}~(ii), (iii), $A^{\mathrm{can}}=A^{\mathrm{act}}\,\Dn_0$ is an identity in $\bbQ(\zeta)[L,L^{-1}]\fps z$, with $A^{\mathrm{can}}_n=L^{4n+1}\mathrm{ph}^{\mathrm{can}}_n(L^{-5})$. Evaluation at $L=1$, that is $u=1$, $q=0$, $\zh=z$, is a ring homomorphism $\bbQ(\zeta)[L,L^{-1}]\to\bbQ(\zeta)$; it gives $A^{\mathrm{act}}(z)|_{L=1}\,\Dn_0(z)=\Phican(z,1)$. As $(1-5^5q)^{-1/5}=1+O(q)$, $A^{\mathrm{act}}_n|_{L=1}$ is the coefficient of $q^0$ of $A^{\mathrm{act}}_n$, which is $h_n$ by (QRR) (\cref{thm:set-qrr}). As $h\in1+z\bbQ\fps z$ is invertible, $\Dn_0=\Phican(z,1)/h(z)$ has rational coefficients.
\end{proof}

\begin{cor}[{\Lr{cor:con-A2}}]\genusstatus{AI-proposed, unchecked}
\label{cor:con-A2}
$\Dn=1$ if and only if $\PhiSP(z,1)=h(z)$. By \cref{thm:con-Pa}~(iv), $\PhiSP(z,1)=c(z;1)$, the value at $L=1$ of the Laurent polynomials $c_n$; so the condition reads $c(z;1)=h(z)$.
\end{cor}

The glue is the place where the value of (QRR) at colour $0$ enters. The coefficient of $z^n$ of $\Dn_0$ enters $\Fg{g}$ first at $g=\lceil(n+3)/3\rceil$ \Lr{rmk:con-genus}.

